\section{\sys: A Zero-Trust Write Boundary}
\label{sec:method}

\sys{} decomposes into the three roles NIST SP 800-207~\cite{nist800207} names.
A Policy Information Point assigns labels and resolves request context; a
Policy Decision Point applies Bell--LaPadula, Biba, and a weighted risk score;
a Policy Enforcement Point acts on the write.

The invariant that names the system: \textbf{labels are recomputed per request
and never cached because data is ``internal'' or ``already in the context
window.''}

\subsection{Architecture}
\label{sec:method:arch}

\subsubsection{Two enforcement planes}
\label{sec:method:planes}

The single most consequential structural decision is that the architecture has
\emph{two} enforcement planes, not one, and they are enforced by different
mechanisms against different objects:

\begin{itemize}
  \item The \textbf{content plane} governs \emph{what text may be emitted}. It
  is the Bell--LaPadula plane, its object is a span, and it is enforced by
  admission, redaction, and abstraction.
  \item The \textbf{action plane} governs \emph{what the agent may cause}. It
  is the Biba plane, its object is a capability, and it is enforced by
  \Cref{eq:downgrade} against the permission system.
\end{itemize}

The separation is forced by the data rather than chosen for symmetry.
\Cref{tab:incidents} contains incidents in which \textbf{no text was emitted at
all}: the PocketOS agent deleted a volume and the Replit agent dropped tables.
A defense that operates only on the content plane---which is every emission
filter, every output guardrail, and every ``do not reveal sensitive
information'' instruction---would have been a no-op in both cases, because
there was no output to filter. Conversely, a purely capability-based defense
cannot express the read-but-never-emit requirement that the high-$\Cl$,
high-$\Il$ quadrant demands, since the issue there is not what the agent may
\emph{do} but what it may \emph{say} while doing it.

This is also why we resist the natural simplification of folding $\Il$ into
$\Cl$ as a single ``risk'' scalar. The two labels feed different planes; a
scalar would have to pick one.

\subsubsection{The request pipeline}
\label{sec:method:pipeline}

A single request traverses six stages. \Cref{tab:stages} states, for each, the
object it operates on, what enforces it, and---the column that matters for a
zero-trust argument---whether a model is in the loop. \todo{Figure 2: render
this pipeline as a block diagram with the two ceilings ($\clr$ and $\lvlI$)
marked and the two planes as parallel tracks. Source in design/pep-design.md.}

\begin{table}[t]
\centering
\caption{The request pipeline. Note the third column: the two stages that do
the most security work involve no model at all, which is what makes them
immune to the compliance failure CI-Work measures.}
\label{tab:stages}
\footnotesize
\begin{tabular}{@{}llcl@{}}
\toprule
\textbf{Stage} & \textbf{Operates on} & \textbf{Model?} & \textbf{Enforced by} \\
\midrule
0. Admission   & connectors, credentials & no  & capability grant, ACL \\
1. PIP context & the request             & no  & directory lookup \\
2. Label       & each entry              & opt & provenance rubric \\
3. PDP         & entry $\times$ request  & no  & \Cref{eq:risk}, precedence \\
4. PEP         & spans and capabilities  & no  & redact / abstract / cap \\
5. Emit        & the artifact            & yes & handle-resolving renderer \\
\bottomrule
\end{tabular}
\end{table}

\paragraph{Stage 0, admission.} Before anything is retrieved, the capability
grant scopes connectors to this one task and the reviewer's clearance bounds
what is readable (precedence rules 1--2, \Cref{sec:method:precedence}). This
stage is cheap, deterministic, and it is where the ambient-credential
incidents would have been stopped: a token that is not in the grant is not
reachable, as opposed to being reachable and later regretted.

\paragraph{Stage 1, context resolution.} The PIP resolves direction, recipient
clearance $\clr(r)$, and required outbound integrity $\lvlI(\afin)$ via
\Cref{eq:lvlI}. \textbf{The two ceilings established here are the whole
design.} In the worked case of \Cref{sec:eval} the agent may \emph{read} at
$0.9$ and may \emph{write} at $0.2$. Every leak lives in that gap, and so does
all of the context that makes the output good.

\paragraph{Stages 2--3, labeling and decision.} Each retrieved entry receives
$\Cl$ and $\Il$; the PDP then applies both rules and the risk score. The
decision is deterministic \emph{given the labels}, which is what makes it
auditable: every reported number traces to a configuration file and a labeled
entry, not to a model's judgement in the moment.

\paragraph{Stages 4--5, enforcement and emission.} The enforcement point acts
on both planes---redacting or abstracting on the content plane, capping the
permitted action on the action plane---and the renderer resolves citation
handles against policy, refusing to expand any handle above the ceiling.

\subsubsection{Where the trust actually sits}

A zero-trust design is only as good as its honesty about residual trust, so we
state it plainly. \sys{} trusts the identity directory, the connector ACL, the
policy configuration, and the offline abstractor's reviewed output. It does
\emph{not} trust the agent, the retrieved context, the recipient, or---and this
is the part that distinguishes it from prompt-level defenses---its own
labeler's content signals over provenance. \Cref{sec:discussion:bugs} describes
what happened the one time we let content override provenance.

\subsection{PIP: the two labels}
\label{sec:method:labels}

\paragraph{Confidentiality $\Cl(e)$.} In \emph{evaluation mode} we take
CI-Work's ground truth ($\text{sensitive}\mapsto 1$,
$\text{essential}\mapsto 0$) so that policy quality is separable from labeler
error---a separation we preserve in every reported result. In \emph{inference
mode} a classifier reusing CI-Work's own nine-category tagging scheme maps to a
$\Cl$ level conditioned on $(\text{sender}, r, d)$.

The recipient clearance $\clr(r)$ is \emph{not} a hyperparameter. It is the
maximum confidentiality the write's audience may receive, which for our setting
is a property of the thread's reader set and is obtained by lookup rather than
judgement. One subtlety carries real weight in deployment: because a thread's
reader set may grow after the write, the correct ceiling is the minimum
clearance over current \emph{and foreseeable} readers, not the current
minimum.

\paragraph{Integrity $\Il(e)$.} This is the new axis, and it must not be a
matter of taste or Tian--Song's critique applies to us as well. We compute it
from a deterministic, auditable provenance rubric over four observable signals:
the source channel (a signed document store outranks official email, which
outranks a chat message, which outranks an author-controlled description);
artifact-type markers (contract, board minute, ledger, audit raise it;
draft, speculation, rumor lower it); author authority; and manipulation
markers, where content proposing deception is forced low.

The governing principle, which \Cref{sec:discussion:bugs} explains we learned
the hard way:
\begin{quote}
\emph{Integrity is a property of provenance, and content cannot vouch for its
own provenance.}
\end{quote}
Concretely, a content marker may \emph{refine} the channel's base score within
a bounded range but may never override it, and author-controlled channels
receive an absolute ceiling applied last so that no later signal can undo it:
\begin{align}
\text{base}   &= \textstyle\max\{\sigma(k) : k \text{ matches } \mathrm{src}(e)\}
                 \;\;\text{or}\;\; \pi_I \text{ if none}\\
\text{score}  &\leftarrow \min\!\big(\max(\text{base}, 0.85),\;
                 \text{base} + \mu\big) \quad \text{[high markers]}\\
\text{score}  &\leftarrow \min(\text{score}, 0.35) \quad \text{[low markers]}\\
\text{score}  &\leftarrow \min(\text{score}, 0.15) \quad \text{[fabrication]}\\
\Il(e)        &= \min(\text{score}, \kappa) \;\;\text{if author-controlled}
\end{align}
with content-uplift bound $\mu = 0.25$, author-controlled ceiling
$\kappa = 0.25$, and prior $\pi_I$. Expressing the bound \emph{relative to the
source's base score} rather than as a per-source ceiling table is what prevents
an unrecognized source from escaping the clamp by failing to match any key.
The asymmetry is intentional and safe: downward markers need no clamp, because
an adversary who lowers the integrity of their own content achieves nothing.

\paragraph{Required outbound integrity $\lvlI(\afin)$.} This composes a floor
from the flow direction with a floor from the agent's authorized capability:
\begin{equation}
\lvlI(\afin) = \max\big(\lvlI^{\mathrm{dir}}(d),\;
                        \lvlI^{\mathrm{cap}}(\mathrm{cap}(A))\big).
\label{eq:lvlI}
\end{equation}
\Cref{eq:lvlI} is where the answer to ``is this task-specific?'' lives, and the
answer is that it is \emph{capability}-specific and only that. The evidence bar
scales with what the write is authorized to \emph{cause}, never with what the
user asked for. Two deployments of the same model reviewing the same pull
request with the same prompt get different integrity requirements if one may
merge and the other may only comment. Nothing about the instruction changes
this; revoking the permission does.

\subsection{PDP: two rules and a weighted score}
\label{sec:method:pdp}

For each candidate entry the two classical rules are evaluated directly:
\begin{align}
\mathrm{ok}_C(e) &:= \Cl(e) \le \clr(r) + \varepsilon_C
  &&\text{(BLP $\star$: no write down)}\\
\mathrm{ok}_I(e) &:= \Il(e) \ge \lvlI(\afin) - \varepsilon_I
  &&\text{(Biba $\star$: no write up)}
\end{align}
yielding a three-valued decision
\begin{equation}
\mathrm{dec}(e) =
\begin{cases}
\textsc{allow} & \mathrm{ok}_C \wedge \mathrm{ok}_I\\
\textsc{deny} & \neg\,\mathrm{ok}_C\\
\textsc{quarantine} & \mathrm{ok}_C \wedge \neg\,\mathrm{ok}_I
\end{cases}
\end{equation}
Confidentiality dominates: dropping an entry protects against leakage
regardless of its integrity, so $\neg\,\mathrm{ok}_C$ short-circuits.
\textsc{quarantine} is the cell that a binary allow/deny policy cannot
represent---an entry safe to mention but not to rely on.

Following Tian--Song, a weighted risk score supplies graduated leniency near
the boundary:
\begin{equation}
\mathrm{risk}(e) = \wC(d)\max\big(0, \Cl(e)-\clr(r)\big)
                 + \wI(d)\max\big(0, \lvlI(\afin)-\Il(e)\big).
\label{eq:risk}
\end{equation}
An entry whose decision is not already \textsc{allow} is released if
$\mathrm{risk}(e) \le \theta_d$ \emph{and} it is not hard-denied.

\paragraph{The initial-trust rule}, which Tian--Song omit and the surveys
criticize them for: every entity starts at a stated prior. Recipient clearance
comes from the lattice; entry integrity comes from the provenance rubric with a
documented default $\pi_I$ when no signal fires; and the agent itself starts at
zero trust, earning nothing by virtue of being the agent. No weight is set
without a stated source.

\subsection{Rule precedence}
\label{sec:method:precedence}

The order of checks is load-bearing rather than cosmetic, and one ordering
constraint is a safety property. \Cref{tab:precedence} gives the six checks.

\begin{table}[t]
\centering
\caption{Admission precedence. Rules 1--3 are enforced before retrieval or
outside the tunable parameters entirely; rule 3 must sit \emph{above} the
$\theta$ override, or regulated content becomes releasable at some operating
point (\Cref{sec:discussion:bugs}).}
\label{tab:precedence}
\footnotesize
\begin{tabular}{@{}clp{0.42\columnwidth}@{}}
\toprule
\# & \textbf{Check} & \textbf{Outcome} \\
\midrule
1 & Need-to-know & Never connected; capability grant excludes the source \\
2 & BLP no-read-up & Not retrievable at this reviewer's clearance \\
3 & Regulated hard block & Denied; \textbf{not eligible for the $\theta$ override} \\
4 & BLP $\star$-property & Read-only: informs reasoning, never emitted \\
5 & Biba no-read-down & Advisory: may raise a hypothesis, never cited as evidence \\
6 & Admitted & Quotable \\
\bottomrule
\end{tabular}
\end{table}

Rules 1 and 2 matter more than their brevity suggests, because they are
enforced with no model in the loop. Under rule 1 the capability grant scopes
connectors to the specific task, so sources with no need---HR systems, vendor
agreements---are not filtered later, they are \emph{never reachable}. This is
least privilege per request, and it is the difference between a scoped grant
and standing OAuth access.

\subsection{PEP: enforcement modes and deployment points}
\label{sec:method:pep}

The enforcement point wraps the final-action tool. It supports pass-through;
\textsc{redact}, removing spans tied to denied entries while keeping the rest;
\textsc{quarantine}, either dropping a low-integrity entry or downgrading it to
attributed form so it is never asserted as fact; and \textsc{deny-and-explain},
a strict mode that upper-bounds privacy and lower-bounds utility.

The \emph{deployment point} turns out to matter more than the mode.
\textbf{Post-hoc} enforcement operates on a finished draft and can therefore
only delete. \textbf{Inline} enforcement runs the PDP \emph{before} generation
and hands the writer a context that already satisfies both rules. We
expected this to be a difference of convenience. \Cref{sec:eval:placement}
shows it is worth all of the conveyance that post-hoc loses, for a structural
reason: two failures are unreachable by deletion. An essential entry above the
write ceiling cannot be paraphrased into an emittable form by a filter, and a
finding the agent dropped because it believed a corrupting entry cannot be
restored, because a filter cannot add text.

\subsection{The enforcement ladder: what makes a decision hold}
\label{sec:method:ladder}

A policy that decides correctly is worthless if the mechanism enforcing it is a
sentence in a system prompt. We therefore separate \emph{what is decided} from
\emph{what makes it hold}, and order the available mechanisms by what an
adversary or a careless model must defeat.

\begin{itemize}
  \item \textbf{L0, prompt instruction.} Enforced by model compliance. The
  model holds the restricted text and chooses. CI-Work has already measured
  this configuration, so it is our baseline row rather than a candidate
  defense: the best prompt intervention leaves violation above 22\% and costs
  8 points of conveyance. Anything enforced by instruction inherits that floor.

  \item \textbf{L1, post-hoc emission filter.} Enforced by a discloser
  detector over the draft. Its achievable violation rate is bounded below by
  the judge's miss rate, and the misses are not random: they cluster on
  paraphrase, on inference from conjunctions of admitted entries (which a
  per-entry filter cannot see by construction), and on structural signals such
  as conspicuously declining to comment on one file. The honest framing is that
  \textbf{L1 bounds careless disclosure, not adversarial disclosure.}

  \item \textbf{L2, handle-mediated citation.} Enforced by the writer not
  possessing the text. At labeling time each above-ceiling entry receives a
  handle and a low-$\Cl$ abstract, itself label-checked against the ceiling.
  Retrieval hands the writer the abstract; the renderer resolves handles
  against policy and refuses to expand any handle above the ceiling. The writer
  can reason (``this defect has precedent here, raise severity''), can cite by
  handle, and cannot inline the specifics because they were never in its
  context. \textbf{This makes no-write-down vacuous rather than filtered for
  the abstracted set.}

  \item \textbf{L3, privileged-context split.} A reasoner sees high-$\Cl$
  entries but can emit only a fixed schema---a verdict enum, a severity level,
  a set of admitted handles to cite, and optionally a bounded reason code---and
  a writer generates prose from admitted entries plus that schema. Restricted
  content influences the verdict through a channel too narrow to carry it. This
  is the strong form, for regulated entries where even an abstract is
  unacceptable, and it is where our design meets the dual-LLM
  pattern~\cite{willison2023dual}; the difference is that the interface
  capacity is declared and bounded rather than left as free text. The residual
  surface is the schema itself: a sufficiently expressive reason-code
  vocabulary becomes a covert channel, and severity levels leak a few bits. We
  treat that as a declared channel capacity rather than pretending it is zero.
\end{itemize}

We implement L1 and a first cut of L2, and report the measured effect of moving
between them in \Cref{sec:eval:placement}.

\paragraph{The trust boundary moves to the abstractor, and we say so.}
Something must read the restricted text to write its abstract. Three properties
make this a smaller problem than the original rather than a relocation of it:
the abstractor runs \emph{offline at labeling time}, once per entry rather than
per request, so its output is reviewable before any agent sees it; its output
is \emph{label-checked} by the same classifier, and an abstract scoring above
the ceiling is rejected and regenerated; and it is a \emph{fixed, small,
auditable corpus} that a human can spot-check in full, which per-request
generations cannot be.

\subsection{Biba as action downgrade, not text filtering}
\label{sec:method:downgrade}

This is the mechanism we consider the paper's most useful single idea, and it
follows from \Cref{eq:bibaviolation} being quantified over evidence rather than
disclosure.

The mistake to avoid is treating integrity as a filtering problem. Attribution
(``an unverified note in the description claims security approved this'') is
genuinely useful, but note precisely what it does: it lowers the \emph{reader's}
reliance. It does not lower the \emph{agent's}. If the verdict already rests on
that entry, attributing it in the prose leaves the wrong verdict standing with
a politer citation.

What enforces Biba for an agent is a capability constraint:
\begin{equation}
\mathrm{act}(A) = \max\Big\{ a \;:\;
  \lvlI(a) \le \min_{e \in \mathrm{ev}(\afin)} \Il(e) \Big\}.
\label{eq:downgrade}
\end{equation}
The integrity of the action the agent may take equals the integrity of the
weakest evidence its conclusion rests on. \Cref{tab:downgrade} instantiates the
ladder for a PR-review agent.

\begin{table}[t]
\centering
\caption{Action downgrade (\Cref{eq:downgrade}): the permitted action follows
from the weakest evidence the verdict rests on. Requires no
natural-language inference and cannot be overridden by context.}
\label{tab:downgrade}
\footnotesize
\begin{tabular}{@{}lp{0.46\columnwidth}@{}}
\toprule
\textbf{Weakest evidence} & \textbf{Highest permitted action} \\
\midrule
$\Il \ge 0.85$ & approve with auto-merge \\
$\Il \ge 0.70$ & approve, human merges \\
$\Il \ge 0.55$ & request changes (blocking) \\
$\Il \ge 0.40$ & non-blocking comment \\
below $0.40$   & ask a question; assert nothing \\
\bottomrule
\end{tabular}
\end{table}

Three properties make this the right primitive. It needs \textbf{no NLP}, being
a comparison over labels already computed. It \textbf{cannot be talked out of},
because it is enforced by the permission system rather than by the model's
cooperation---which makes it immune to the 15.8--50.9\% compliance failure
CI-Work measures. And it \textbf{degrades gracefully}: weak evidence yields a
weaker action rather than a refusal, so utility falls off smoothly instead of
collapsing the way naive denial does.

Both mechanisms ship together: attribution for reader-facing honesty, action
downgrade for correctness.

\subsection{Deriving the weights}
\label{sec:method:weights}

A reviewer who accepts our critique of Tian--Song's ad-hoc weights will
immediately ask where ours come from. We answer by splitting every parameter
into three tiers with \emph{different justification obligations}, under one
governing principle:
\begin{quote}
\emph{Anything encoding a norm must be derived from outside the benchmark.
Anything derived from the benchmark is a measurement parameter and must be
reported with a held-out estimate.}
\end{quote}
Fitting recipient clearance to the benchmark would define away the thing being
measured, since ``sensitive'' would become whatever the policy blocks.

\paragraph{Tier 1, structural: never fit.} $\clr(r)$, $\lvlI$, and the
regulated-class set come from the deployment's existing access control and the
agent's authorized action surface. In deployment $\clr(r)$ is replaced by a
query against repository visibility and team membership; that substitution is
what makes the design deployable and is the reason this tier must not be
tuned. We note explicitly that $\clr$ is \emph{not} $\wC$: the Upward
direction has the highest clearance, because a superior may receive almost
anything, while also being the direction CI-Work found leaks most. Clearance is
what is \emph{permitted}; the observed leak rate is how often models
\emph{fail}. Conflating them would be a category error.

\paragraph{Tier 2, measurement: calibrate and report error.} The provenance
rubric's weights are fit for agreement with human integrity labels on a
development split and then frozen, following CI-Work's own annotation protocol
so the numbers are comparable: two annotators label a dev sample on a four-point
scale blind to the rubric, weights are fit on dev only, and agreement is
reported on a held-out split. CI-Work reports 82.5--95.0\% for entry labels and
83.0--91.0\% for its judge; materially below that band is not publishable as a
labeler. \needsrun{this protocol has not been executed; $\pi_I = 0.5$ is
currently a placeholder chosen for symmetry, which is not a justification}

\paragraph{Tier 3, operating point: derive, then report a curve.} The insight
that removes the hand-tuning is that $\wC$ and $\wI$ are not free parameters.
They are the relative cost of a unit breach on each axis, and cost is something
a deployment can state. We derive $\wC$ from audience breadth and $\wI$ from
consequence class, each normalized to the Lateral reference
(\Cref{tab:weights}):
\begin{equation}
\wC(d) = \frac{\mathrm{breadth}(d)}{\mathrm{breadth}(\text{Lateral})},
\qquad
\wI(d) = \frac{\mathrm{consequence}(d)}{\mathrm{consequence}(\text{Lateral})}.
\end{equation}

\begin{table}[t]
\centering
\caption{Derived operating-point weights. $\clr$ is Tier 1 (an ACL lookup);
$\wC$ and $\wI$ are Tier 3, derived from audience breadth and consequence class
rather than tuned. Note $\clr$ and $\wC$ move in \emph{opposite} directions:
Upward permits the most and costs the least per breach.}
\label{tab:weights}
\footnotesize
\begin{tabular}{@{}lccccc@{}}
\toprule
\textbf{Direction} & \textbf{Audience} & $\mathrm{breadth}$ & $\wC$ & $\wI$ & $\clr$ \\
\midrule
Downward & own team      & 2  & 0.67 & 0.6 & 0.7 \\
Upward   & own team      & 2  & 0.67 & 0.6 & 0.9 \\
Lateral  & two teams     & 3  & 1.00 & 1.0 & 0.6 \\
Diagonal & org-wide      & 4  & 1.33 & 1.6 & 0.4 \\
External & public        & 10 & 3.33 & 2.0 & 0.2 \\
\bottomrule
\end{tabular}
\end{table}

\paragraph{A calibration hazard we must flag.} Because \Cref{eq:risk} is
denominated in weighted units, rescaling any weight silently changes what
$\theta$ means. Moving $\wC(\text{External})$ from $1.0$ to $3.33$ makes the
same $\theta=0.15$ more than three times stricter on the confidentiality axis.
Two consequences are mandatory rather than optional: $\theta$ must be re-swept
whenever any weight changes and must never be inherited across weight
configurations, and results must report the \emph{curve} over $\theta$ together
with the chosen operating point and its cost justification. A single $\theta$
with no curve is not defensible. \needsrun{the $\theta$ sweep}

\subsection{Graduated severity}
\label{sec:method:severity}

CI-Work's Leakage is a mean and its Violation is a max, so neither can
represent the difference between disclosing an internal ticket identifier and
disclosing a named customer from a postmortem. We add a multiplicative
per-entry severity:
\begin{equation}
\mathrm{sev}(e,d,k) = \mathrm{base}_C(e)\cdot \mathrm{breadth}(d)
  \cdot \mathrm{persist}(k) \cdot \mathrm{reg}(e)
\label{eq:severity}
\end{equation}
with $\mathrm{base}_C \in \{0,1,4,16\}$ over Public/Internal/Confidential/
Restricted, $\mathrm{breadth} \in \{1,2,4,10\}$, $\mathrm{persist} \in \{1,3\}$
for ephemeral chat versus an indexed thread, and
$\mathrm{reg} \in \{1,3,5,8\}$ for none/PII/PCI-HIPAA/export-controlled.
$\mathrm{base}_C$ is geometric deliberately: on a linear scale a dozen Internal
disclosures would outweigh one Restricted disclosure, inverting the ordering
any real organization wants. The resulting spread between the cheapest and
costliest disclosure in a single thread is $240\times$, which a mean averages
away and a max cannot see.
